\documentclass[12pt,a4paper]{ctexart}
\usepackage[margin=2.5cm]{geometry}
\usepackage{array}
\usepackage{booktabs}
\usepackage{tabularx}

\renewcommand{\arraystretch}{1.5}
\pagestyle{plain}

\begin{document}

\begin{center}
	{\zihao{2}\bfseries AI工具使用详情}
\end{center}

\noindent\textbf{填写说明：}下表仅为排版示例。提交前应删除本说明，并将全部内容替换为本队真实的
AI 工具使用记录；不得编造或替换工具名称、版本及使用过程。

\begin{table}[htbp]
	\centering
	\caption{AI工具使用详情示例}
	\begin{tabularx}{\textwidth}{>{\bfseries\raggedright\arraybackslash}p{3.6cm}X}
		\toprule
		记录项目 & 示例内容（请按实际情况修改） \\
		\midrule
		工具名称、版本或型号 & DeepSeek，DeepSeek-R1-0528（仅为国产模型排版示例） \\
		具体使用目的和环节 & 在模型求解阶段辅助定位 Python 数组维度错误，并检查绘图代码中的坐标轴标签。 \\
		主要提示方式与使用过程 & 提供报错信息、相关函数及预期输出，请工具分析错误原因并给出排查步骤；团队逐项复现后再修改代码。 \\
		采纳、人工修改和核验 & 仅采纳数组维度检查思路，代码由团队重新编写，并使用原始数据重复运行、对照关键结果完成核验。 \\
		\bottomrule
	\end{tabularx}
\end{table}

\end{document}
